\subsection{Réactivité du lithium métallique}
\label{subsection:reac_Li_m}
\setcounter{equation}{0}
\setcounter{Q}{0}

\begin{figure}[H]
\begin{centering}
\subfloat[\label{fig:05-01}]{
\includegraphics[height=3cm]{Ex05/Fig01}}
\hfill
\subfloat[\label{fig:05-02}]{\includegraphics[height=3cm]{Ex05/Fig02}}
\hfill
\subfloat[\label{fig:05-03}]{\includegraphics[height=3cm]{Ex05/Fig03}}
\end{centering}
\caption{Quelques phases courantes du lithium : \ref{fig:05-01} le lithium métal, \ref{fig:05-02} l'oxyde de lithium, \ref{fig:05-03} le nitrure de lithium.}
\end{figure}

Le métal lithium réagit spontanément avec l'eau lors d'une réaction fortement exothermique, relativement lente, libérant un ou plusieurs gaz. On observe expérimentalement un abaissement du pH, appelé alcalinisation du milieu.

\Q Donner l'équation de réaction entre le couple \ce{Li+}/\ce{Li} et l'un des couples de l'eau.
\solution{Le lithium est un métal très réducteur, il permet la réduction de l'eau.  Associons les couples \ce{H^+ (aq) / H2 (g)} et \ce{Li^+ (aq) / Li (s)}.
\begin{equation*}
\ce{ Li (s) + H^+ (aq) = Li^+ (aq) + \frac{1}{2} H2 (g)}
\end{equation*}
Le couple \ce{H2O (aq) / H2 (g)} peut aussi être utilisé comme c'est le cas à l'exercice suivant.
}


En présence d'atmosphère, que l'on supposera constituée uniquement de \ce{O2} et \ce{N2}, le lithium conduit à la formation de nitrures de lithium et d'oxydes de lithium. 

\Q Après avoir rappelé la configuration électronique des atomes de \ce{Li}, \ce{N} et \ce{O} en déduire la formule chimique du nitrure et de l'oxyde de lithium.
\solution{
D'après les règles de l'Aufbau, on a :
\begin{align*}
_3 [\ce{Li}] &= _2[\ce{He}] 2s^1 \\
_7 [\ce{N}] &= _2[\ce{He}] 2s^2 2p^3 \\
_8 [\ce{O}] &= _2[\ce{He}] 2s^2 2p^4 \\
\end{align*}
Les données nous indiquent que $\chi ( \ce{O}) > \chi( \ce{Li} )$ et $\chi ( \ce{N} ) > \chi ( \ce{Li} )$, par conséquent, dans ces structures, le lithium est oxydé et l'azote/oxygène sont réduits. Pour se ramener à la règle de l'octet, on s'attend donc à avoir $no(\ce{Li})=+I$, $no(\ce{N}) = -III$ et $no(\ce{O}) = -II$
\begin{align*}
\ce{Li_x N} \implies x \cdot no (\ce{Li}) + no(\ce{N}) = 0 \\
\ce{Li_y O} \implies y \cdot no (\ce{Li}) + no(\ce{O}) = 0
\end{align*}
Par conséquent, les formules brutes sont \ce{Li3N} pour le nitrure de lithium et \ce{Li2O} pour l'oxyde de lithium.}

\Q Exprimer les demi-équations rédox et l'équation de réaction pour le nitrure et pour l'oxyde.
\solution{Pour le nitrure, on définit deux couples rédox : \ce{Li^+ / Li (s)} et \ce{N2 (g)/ Li3N (s)}.
\begin{center}
\begin{minipage}{0.4\textwidth}
\begin{align*}
\ce{Li(s) &= Li^+ + e^-} \quad \times 3 \\
\ce{\frac{1}{2} N2 (g) &+ 3 Li^+ + 3 e^- = Li_3 N (s)} \\
\hline 
\ce{\frac{1}{2} N2 (g) &+ 3 Li(s) = Li_3 N (s)}
\end{align*}
\end{minipage}
\end{center}
Pour l'oxyde, on définit deux couples rédox : \ce{Li^+ / Li (s)} et \ce{O2 (g)/ Li2O (s)}.
\begin{center}
\begin{minipage}{0.4\textwidth}
\begin{align*}
\ce{Li(s) &= Li^+ + e^-} \quad \times 2 \\
\ce{\frac{1}{2} O2 (g) &+ 2 Li^+ + 2 e^- = Li_2 N (s)} \\
\hline 
\ce{\frac{1}{2} O2 (g) &+ 3 Li(s) = Li_2 O (s)}
\end{align*}
\end{minipage}
\end{center}}


%Le nitrure de lithium est un composé de couleur noir et isolant.
\Q Comment peut-on remarquer "à l'{\oe}il" la présence de nitrures de lithium ? Quelles précautions doit-on prendre lorsque l'on manipule du lithium métallique ?
\solution{Le nitrure de lithium apparaît spontanément (au sens thermodynamique) dès la mise en contact du lithium métallique avec l'air, lequel prend une teinte noire en surface. \\
Afin d'éviter qu'il y ait contamination par l'air, les synthèses nécessitant du lithium métallique ou dont l'un des intermédiaires de synthèse est le lithium métallique doivent être réalisées en boîte à gant pressurisé rempli d'argon ou d'hélium.
\begin{center}
\includegraphics[width=0.4\textwidth]{Ex05/Corr01.png}
\end{center}
La présence de nitrures de lithium peut être fortement dommageable sur les propriétés des matériaux (\emph{par ex.} le nitrure de lithium est un isolant électrique) plutôt cassant, alors que le lithium métallique est un métal ductile et conducteur électrique.}

\paragraph*{Données} : 
\begin{itemize}
\item Numéro atomique : $Z(\ce{Li}) = 3$, $Z(\ce{N}) = 7$, $Z(\ce{O}) = 8$  
\item \'{E}lectronégativité de Pauling : $\chi( \ce{Li} ) = 2.1$ , $\chi( \ce{N} ) = 3.04$ et $\chi( \ce{O} ) = 3.44$
\end{itemize}